% ============================================================
%  Магистерская диссертация — TechStore
%  Компилятор: XeLaTeX (в Overleaf: Menu → Compiler → XeLaTeX)
%  Команда:    xelatex magisterskaya_dissertatsiya.tex  (дважды)
% ============================================================
\documentclass[14pt,a4paper]{extreport}

% ── Шрифты и язык (XeLaTeX) ─────────────────────────────────
\usepackage{fontspec}
\setmainfont{Times New Roman}
\setsansfont{Arial}
\setmonofont[Scale=0.85]{Courier New}

\usepackage{polyglossia}
\setmainlanguage{russian}
\setotherlanguage{english}

% ── Поля страницы (ГОСТ 7.32) ───────────────────────────────
\usepackage[left=30mm,right=15mm,top=20mm,bottom=20mm]{geometry}

% ── Межстрочный интервал 1.5 ────────────────────────────────
\usepackage{setspace}
\onehalfspacing

% ── Абзацный отступ ─────────────────────────────────────────
\usepackage{indentfirst}
\setlength{\parindent}{1.25cm}

% ── Таблицы ─────────────────────────────────────────────────
\usepackage{booktabs}
\usepackage{longtable}
\usepackage{array}

% ── Листинги кода ───────────────────────────────────────────
\usepackage{listings}
\usepackage{xcolor}
\lstset{
  basicstyle=\small\ttfamily,
  breaklines=true,
  frame=single,
  numbers=left,
  numberstyle=\tiny\color{gray},
  keywordstyle=\color{blue!70!black},
  commentstyle=\color{gray},
  stringstyle=\color{red!60!black},
  showstringspaces=false,
  captionpos=b,
  tabsize=2,
}

% ── Математика ──────────────────────────────────────────────
\usepackage{amsmath}

% ── Гиперссылки ─────────────────────────────────────────────
\usepackage[hidelinks,unicode]{hyperref}

% ── Заголовки глав ──────────────────────────────────────────
\usepackage{titlesec}
\titleformat{\chapter}[block]
  {\normalfont\large\bfseries\centering}
  {\thechapter.}{0.5em}{}
\titlespacing*{\chapter}{0pt}{12pt}{12pt}
\titleformat{\section}[block]
  {\normalfont\normalsize\bfseries}
  {\thesection}{0.5em}{}
\titleformat{\subsection}[block]
  {\normalfont\normalsize\bfseries}
  {\thesubsection}{0.5em}{}

% ── Колонтитулы ─────────────────────────────────────────────
\usepackage{fancyhdr}
\pagestyle{fancy}
\fancyhf{}
\fancyhead[R]{\small\thepage}
\renewcommand{\headrulewidth}{0pt}

% ============================================================
\begin{document}

% ── Титульный лист ───────────────────────────────────────────
\begin{titlepage}
\begin{center}
{\small МИНИСТЕРСТВО НАУКИ И ВЫСШЕГО ОБРАЗОВАНИЯ РЕСПУБЛИКИ КАЗАХСТАН}

\vspace{0.5cm}
{\small \textbf{[Название университета]}}

\vspace{0.3cm}
{\small Кафедра информационных систем и технологий}

\vspace{2cm}
{\Large \textbf{МАГИСТЕРСКАЯ ДИССЕРТАЦИЯ}}

\vspace{1cm}
{\normalsize \textbf{Тема:}}

\vspace{0.5cm}
{\large \textbf{Проектирование и разработка многофункциональной информационной системы электронной торговли с поддержкой программы лояльности и многоязычного интерфейса}}

\vspace{2cm}
\begin{flushleft}
\hspace{6cm}
\begin{minipage}{0.5\textwidth}
\textbf{Специальность:} 7М06101 --- Информационные системы

\vspace{0.5cm}
\textbf{Выполнил:} [ФИО магистранта]

\vspace{0.5cm}
\textbf{Научный руководитель:} [ФИО, учёная степень, должность]
\end{minipage}
\end{flushleft}

\vfill
{\normalsize Алматы / Астана --- 2026}
\end{center}
\end{titlepage}

% ── Аннотация ────────────────────────────────────────────────
\chapter*{АННОТАЦИЯ}
\addcontentsline{toc}{chapter}{АННОТАЦИЯ}

Настоящая диссертация посвящена проектированию и практической реализации многофункциональной информационной системы электронной торговли «TechStore», ориентированной на рынок Республики Казахстан. В работе рассматриваются архитектурные решения на базе модульного подхода с использованием стека технологий React, Vite, TailwindCSS (клиентская часть) и FastAPI, SQLAlchemy, PostgreSQL (серверная часть). Реализованы модули: аутентификация с верификацией почты, каталог товаров с нечётким поиском, корзина покупок с тремя способами оплаты, девятиуровневая программа лояльности, многоязычный интерфейс (русский, казахский, английский), модуль доставки по 55 городам Казахстана, административная панель, портал сотрудников и управление складскими остатками по пяти филиалам. Система развёрнута в облачной инфраструктуре Vercel (фронтенд) и Render (бэкенд + PostgreSQL).

\textbf{Ключевые слова:} информационная система, электронная торговля, REST API, программа лояльности, JWT-аутентификация, многоязычный интерфейс, облачный деплой, FastAPI, React.

\chapter*{ABSTRACT}
\addcontentsline{toc}{chapter}{ABSTRACT}

This dissertation is devoted to the design and practical implementation of a multi-functional e-commerce information system ``TechStore'' targeting the market of the Republic of Kazakhstan. The work examines architectural solutions based on a modular approach using a technology stack of React, Vite, TailwindCSS (client side) and FastAPI, SQLAlchemy, PostgreSQL (server side). Modules implemented: authentication with email verification, product catalog with fuzzy search, shopping cart with three payment methods, a nine-tier loyalty program, multilingual interface (Russian, Kazakh, English), delivery module covering 55 cities of Kazakhstan, admin panel, employee portal, and stock management across five branches. The system is deployed on Vercel (frontend) and Render (backend + PostgreSQL).

\textbf{Keywords:} information system, e-commerce, REST API, loyalty program, JWT authentication, multilingual interface, cloud deployment, FastAPI, React.

% ── Оглавление ───────────────────────────────────────────────
\tableofcontents

% ============================================================
\chapter{ВВЕДЕНИЕ}

\section{Актуальность темы исследования}

Электронная коммерция является одним из наиболее динамично развивающихся секторов цифровой экономики. По данным Международной организации электронной коммерции, глобальный объём онлайн-торговли ежегодно прирастает на 15--20\%. Казахстан демонстрирует аналогичную тенденцию: по информации Агентства по статистике РК, объём рынка электронной торговли страны превысил 2 триллиона тенге, а число интернет-покупателей ежегодно увеличивается на 25--30\%.

Несмотря на рост рынка, значительная часть отечественных торговых предприятий либо использует зарубежные платформы (Kaspi, AliExpress), либо располагает устаревшими информационными системами, не адаптированными к локальным особенностям: трёхъязычному интерфейсу (казахский, русский, английский), валюте тенге, региональной логистике и системам рассрочки платежей. Проектирование собственной информационной системы позволяет устранить эти недостатки и создать конкурентоспособный продукт.

\section{Цель и задачи исследования}

\textbf{Цель исследования:} разработать и программно реализовать многофункциональную информационную систему электронной торговли, адаптированную к условиям казахстанского рынка.

\textbf{Задачи исследования:}
\begin{enumerate}
  \item Провести анализ существующих платформ электронной коммерции и выявить их ограничения применительно к казахстанскому рынку.
  \item Разработать концептуальную и логическую модели данных системы.
  \item Спроектировать RESTful API-архитектуру серверной части с использованием FastAPI.
  \item Реализовать клиентскую часть на базе React с поддержкой трёх языков.
  \item Разработать модуль аутентификации с JWT-токенами и верификацией электронной почты.
  \item Создать многоуровневую программу лояльности с начислением кешбэка.
  \item Реализовать модуль доставки с картой 55 городов Казахстана и привязкой к филиалам.
  \item Развернуть систему в облачной инфраструктуре и обеспечить её доступность.
  \item Провести функциональное тестирование всех модулей системы.
\end{enumerate}

\section{Объект и предмет исследования}

\textbf{Объект исследования:} информационная система управления торговыми операциями в сфере электронной коммерции.

\textbf{Предмет исследования:} методы проектирования, архитектурные паттерны и программные средства для построения масштабируемых веб-приложений электронной торговли.

\section{Методы исследования}

В ходе работы применялись следующие методы:
\begin{itemize}
  \item системный анализ предметной области электронной торговли;
  \item сравнительный анализ существующих программных решений;
  \item объектно-ориентированное проектирование (UML-диаграммы);
  \item итеративная разработка (Agile/Scrum);
  \item функциональное и интеграционное тестирование.
\end{itemize}

\section{Научная новизна}

Научная новизна работы заключается в:
\begin{itemize}
  \item разработке комплексной модели программы лояльности с девятью градациями уровней применительно к казахстанскому рынку;
  \item создании модели региональной доставки, основанной на геолокационном сопоставлении 55 городов РК с ближайшими из пяти операционных филиалов;
  \item реализации архитектурного решения для многоязычного интерфейса (ru/kz/en) с динамическим переключением без перезагрузки страницы.
\end{itemize}

\section{Практическая значимость}

Разработанная система может быть непосредственно использована организациями розничной торговли Казахстана как готовое решение с открытым исходным кодом. Архитектурные решения и подходы применимы для разработки аналогичных систем в смежных отраслях.

\section{Структура диссертации}

Диссертация состоит из введения, семи глав, заключения, списка литературы (40 источников) и приложений. Общий объём --- 80--100 страниц.

% ============================================================
\chapter{ОБЗОР ЛИТЕРАТУРЫ}

\section{Электронная коммерция: определения и классификация}

Электронная коммерция (e-commerce) определяется как форма торговых отношений, при которой заключение договоров и операционная деятельность осуществляются с использованием информационно-коммуникационных технологий \cite{turban2018}. В научной литературе выделяют несколько классификационных групп:

\begin{itemize}
  \item по типу участников: B2B (бизнес--бизнес), B2C (бизнес--потребитель), C2C (потребитель--потребитель);
  \item по типу товаров: физические товары, цифровые товары, услуги;
  \item по модели доставки: собственная логистика, агрегаторы доставки, цифровая дистрибуция.
\end{itemize}

Разрабатываемая система TechStore относится к модели B2C с поддержкой как физических (гаджеты), так и цифровых товаров.

\section{Обзор существующих платформ}

\subsection{Зарубежные платформы}

\textit{Shopify} --- облачная SaaS-платформа. Недостатки для казахстанского рынка: отсутствие поддержки казахского языка, ограниченная интеграция с местными платёжными системами, высокая стоимость (от \$39/мес.) \cite{shopify2024}.

\textit{WooCommerce} --- расширение WordPress. Гибкость настройки компенсируется сложностью поддержки и уязвимостями безопасности \cite{woocommerce2024}.

\textit{Magento (Adobe Commerce)} --- энтерпрайз-платформа с высокими затратами на внедрение и лицензирование \cite{adobe2024}.

\subsection{Казахстанские платформы}

\textit{Kaspi.kz} --- лидер местного рынка, закрытая экосистема с комиссией 10--20\%, отсутствием возможности создания брендированного магазина \cite{kaspi2024}.

\textbf{Вывод:} ни одна из рассмотренных платформ не обеспечивает полного контроля над продуктом, гибкой программы лояльности, трёхъязычного интерфейса и интеграции с региональной логистикой при низкой стоимости владения.

\section{Технологический стек: обзор литературы}

\textbf{React.} Библиотека реализует концепцию виртуального DOM для эффективного обновления интерфейса \cite{react2024}. Исследования показывают, что приложения на React демонстрируют на 40\% меньшее время до интерактивности по сравнению с традиционными MPA \cite{garrett2022}.

\textbf{FastAPI.} Фреймворк обеспечивает автоматическую генерацию документации OpenAPI/Swagger, валидацию данных через Pydantic и асинхронную обработку запросов. По бенчмаркам TechEmpower, FastAPI превосходит Flask в пропускной способности на 300\% \cite{fastapi2024}.

\textbf{PostgreSQL.} ACID-совместимость, поддержка JSON-полей, расширенные типы индексов делают её оптимальным выбором для e-commerce систем \cite{postgresql2024}.

\textbf{JWT.} Стандарт RFC 7519 обеспечивает безсостоятельную (stateless) аутентификацию, критичную для горизонтального масштабирования \cite{rfc7519}.

\section{Программы лояльности}

Исследования Reichheld \& Schefter показали, что увеличение удержания клиентов на 5\% повышает прибыль на 25--95\% \cite{reichheld2000}. В системе реализована гибридная модель: статус определяется накопленной суммой покупок, кешбэк зависит от статуса (0\%--50\%). Кешбэк начисляется только при оплате картой.

\section{Интернационализация и облачный деплой}

Казахстан имеет государственный (казахский) и официальный (русский) языки \cite{stat2023}. Использована библиотека i18next с React-адаптером. Модель PaaS на Vercel и Render обеспечивает CI/CD, масштабирование и SSL/TLS \cite{12factor}.

% ============================================================
\chapter{ТЕОРЕТИЧЕСКАЯ ЧАСТЬ}

\section{Архитектурные паттерны}

\subsection{Монолитная vs. Микросервисная архитектура}

Применена \textbf{модульная монолитная архитектура}: единый процесс FastAPI с чёткой декомпозицией роутеров (\texttt{/api/auth}, \texttt{/api/products}, \texttt{/api/orders} и др.). Это обеспечивает простоту деплоя при сохранении логической изоляции модулей.

\subsection{REST API}

HTTP-методы маппируются на CRUD-операции согласно принципам REST \cite{fielding2000}:

\begin{table}[h!]
\centering
\caption{Маппинг HTTP-методов}
\begin{tabular}{lll}
\toprule
\textbf{Метод} & \textbf{Операция} & \textbf{Пример} \\
\midrule
GET    & Чтение     & \texttt{GET /api/products} \\
POST   & Создание   & \texttt{POST /api/orders} \\
PUT    & Обновление & \texttt{PUT /api/products/\{id\}} \\
DELETE & Удаление   & \texttt{DELETE /api/admin/users/\{id\}} \\
\bottomrule
\end{tabular}
\end{table}

\section{Модели данных}

\subsection{Концептуальная ER-модель}

\begin{lstlisting}[language={},caption={Основные связи между сущностями}]
User (1) --< Order (M)
Order (1) --< OrderItem (M)
OrderItem (M) >-- Product (1)
Branch (1) --< BranchStock (M)
BranchStock (M) >-- Product (1)
User (M) >-- Branch (1)   [сотрудник привязан к филиалу]
Order (1) --< CashbackTransaction (M)
\end{lstlisting}

\subsection{Логическая модель}

\begin{table}[h!]
\centering
\caption{Таблица \texttt{users}}
\begin{tabular}{lll}
\toprule
\textbf{Поле} & \textbf{Тип} & \textbf{Описание} \\
\midrule
id               & INTEGER PK       & Первичный ключ \\
name             & VARCHAR(255)     & Имя пользователя \\
email            & VARCHAR(255) UQ  & Электронная почта \\
hashed\_password & VARCHAR(255)     & Хэш пароля (bcrypt) \\
is\_admin        & BOOLEAN          & Признак администратора \\
is\_employee     & BOOLEAN          & Признак сотрудника \\
branch\_id       & FK → branches    & Привязка к филиалу \\
cashback\_balance & FLOAT           & Баланс кешбэка (тг) \\
total\_spent     & FLOAT            & Накопленная сумма покупок \\
email\_verified  & BOOLEAN          & Верификация email \\
created\_at      & DATETIME         & Дата регистрации \\
\bottomrule
\end{tabular}
\end{table}

\begin{table}[h!]
\centering
\caption{Таблица \texttt{orders}}
\begin{tabular}{lll}
\toprule
\textbf{Поле} & \textbf{Тип} & \textbf{Описание} \\
\midrule
id                  & INTEGER PK    & Первичный ключ \\
user\_id            & FK → users    & Покупатель \\
branch\_id          & FK → branches & Филиал обработки \\
delivery\_city      & VARCHAR(100)  & Город доставки \\
total               & FLOAT         & Итоговая сумма (тг) \\
cashback\_used      & FLOAT         & Использованный кешбэк \\
cashback\_earned    & FLOAT         & Начисленный кешбэк \\
payment\_method     & VARCHAR(20)   & card / cash / installment \\
installment\_months & INTEGER       & Число месяцев рассрочки \\
status              & VARCHAR(50)   & pending / processing / completed \\
\bottomrule
\end{tabular}
\end{table}

\section{Программа лояльности: математическая модель}

\begin{table}[h!]
\centering
\caption{Уровни программы лояльности}
\begin{tabular}{clrr}
\toprule
\textbf{№} & \textbf{Уровень} & \textbf{Порог (тг)} & \textbf{Кешбэк} \\
\midrule
1 & Новичок   &           0 &  0\% \\
2 & Бронза    &      50\,000 &  3\% \\
3 & Серебро   &     150\,000 &  5\% \\
4 & Золото    &     300\,000 &  7\% \\
5 & Платина   &     600\,000 & 10\% \\
6 & Сапфир    &   1\,000\,000 & 13\% \\
7 & Изумруд   &   2\,000\,000 & 17\% \\
8 & Бриллиант &   5\,000\,000 & 25\% \\
9 & Легенда   &  10\,000\,000 & 50\% \\
\bottomrule
\end{tabular}
\end{table}

Формула начисления кешбэка при оплате картой:
\begin{equation}
  \text{cashback\_earned} = \text{order\_total} \times \text{tier\_rate}
\end{equation}

Условие: \texttt{payment\_method == "card"}. При наличной оплате и рассрочке кешбэк не начисляется.

\section{Алгоритм доставки}

Для 55 городов Казахстана реализована карта сопоставления с ближайшим из пяти филиалов. Поиск выполняется за O(1) через Python-словарь по названию города.

\section{Аутентификация и безопасность}

\subsection{JWT-поток}
\begin{enumerate}
  \item \texttt{POST /api/auth/token} с email + паролем.
  \item Сервер верифицирует пароль через \texttt{bcrypt.checkpw()}.
  \item Генерируется JWT: \texttt{\{sub: email, exp: now+30d\}}, подпись HS256.
  \item Клиент сохраняет токен, передаёт в заголовке \texttt{Authorization: Bearer}.
\end{enumerate}

\subsection{Верификация email}

UUID-токен генерируется при регистрации, отправляется через Resend API. Срок действия --- 24 часа. Неверифицированные пользователи не могут оформлять заказы.

% ============================================================
\chapter{ПРОЕКТИРОВАНИЕ СИСТЕМЫ}

\section{Функциональные требования}

\textbf{Аутентификация (AUTH):}
FR-01: Регистрация с верификацией email;
FR-02: Вход по email/паролю с выдачей JWT;
FR-03: Сброс пароля через email;
FR-04: Ролевое разграничение (покупатель / сотрудник / администратор).

\textbf{Каталог (CATALOG):}
FR-05: Отображение 30+ товаров;
FR-06: Фильтрация по категории и цене;
FR-07: Нечёткий текстовый поиск;
FR-08: Детальная страница товара.

\textbf{Заказы (ORDERS):}
FR-09: Корзина товаров;
FR-10: Три способа оплаты (карта, наличные, рассрочка 3/6/12 мес.);
FR-11: Выбор города доставки с расчётом срока;
FR-12: История заказов.

\textbf{Лояльность (LOYALTY):}
FR-13: Начисление кешбэка при оплате картой;
FR-14: Отображение уровня и баланса;
FR-15: Использование кешбэка при покупке.

\textbf{Интерфейс (I18N):}
FR-16: Три языка (RU/KZ/EN);
FR-17: Сохранение языка между сессиями.

\textbf{Администрирование (ADMIN):}
FR-19: Дашборд; FR-20: CRUD товаров; FR-21: Заказы; FR-22: Пользователи.

\textbf{Сотрудник (EMPLOYEE):}
FR-23: Заказы филиала; FR-24: Изменение статуса заказа.

\section{Нефункциональные требования}

\begin{table}[h!]
\centering
\caption{Нефункциональные требования}
\begin{tabular}{lp{9cm}}
\toprule
\textbf{Код} & \textbf{Требование} \\
\midrule
NFR-01 & Время ответа API $\leq$ 500 мс (95-й перцентиль) \\
NFR-02 & Доступность $\geq$ 99.5\% \\
NFR-03 & bcrypt, не менее 12 раундов \\
NFR-04 & TLS 1.2+ на всех соединениях \\
NFR-05 & JWT: срок $\leq$ 30 дней, HS256 \\
NFR-06 & CORS: только доверенные домены \\
NFR-07 & Горизонтальная масштабируемость \\
\bottomrule
\end{tabular}
\end{table}

\section{Архитектура системы}

Трёхзвенная архитектура:
\begin{enumerate}
  \item \textbf{Представление:} React 18 SPA на Vercel CDN
  \item \textbf{Бизнес-логика:} FastAPI (Python 3.11) на Render
  \item \textbf{Данные:} PostgreSQL 15 на Render Managed Database
\end{enumerate}

Взаимодействие --- HTTPS/REST API с JWT-аутентификацией.

% ============================================================
\chapter{ПРАКТИЧЕСКАЯ РЕАЛИЗАЦИЯ}

\section{Стек технологий}

\begin{table}[h!]
\centering
\caption{Технологии фронтенда}
\begin{tabular}{lll}
\toprule
\textbf{Технология} & \textbf{Версия} & \textbf{Обоснование} \\
\midrule
React        & 18.x   & Virtual DOM, хуки, Context API \\
Vite         & 5.x    & HMR < 50 мс, оптимальная prod-сборка \\
TailwindCSS  & 3.x    & Utility-first, быстрая вёрстка \\
React Router & 6.x    & Декларативная маршрутизация \\
i18next      & latest & Промышленный стандарт i18n \\
Axios        & latest & HTTP-клиент с интерцепторами \\
\bottomrule
\end{tabular}
\end{table}

\begin{table}[h!]
\centering
\caption{Технологии бэкенда}
\begin{tabular}{lll}
\toprule
\textbf{Технология} & \textbf{Версия} & \textbf{Обоснование} \\
\midrule
Python      & 3.11    & Аннотации типов, богатая экосистема \\
FastAPI     & 0.110+  & Асинхронный, автодокументация OpenAPI \\
SQLAlchemy  & 2.x     & Современный ORM \\
PostgreSQL  & 15      & ACID, JSON-поля, масштабируемость \\
bcrypt      & latest  & Адаптивное хэширование паролей \\
python-jose & latest  & JWT HS256 \\
Resend      & HTTP    & Транзакционные письма \\
\bottomrule
\end{tabular}
\end{table}

\section{Реализация ключевых модулей}

\subsection{Аутентификация}

\begin{lstlisting}[language=Python, caption={Хэширование и проверка пароля}]
import bcrypt

def hash_password(password: str) -> str:
    return bcrypt.hashpw(
        password.encode(), bcrypt.gensalt(rounds=12)
    ).decode()

def verify_password(password: str, hashed: str) -> bool:
    return bcrypt.checkpw(password.encode(), hashed.encode())
\end{lstlisting}

\begin{lstlisting}[language=Python, caption={Генерация JWT-токена}]
from jose import jwt
from datetime import datetime, timedelta

def create_access_token(data: dict) -> str:
    payload = {
        **data,
        "exp": datetime.utcnow() + timedelta(days=30)
    }
    return jwt.encode(payload, SECRET_KEY, algorithm="HS256")
\end{lstlisting}

\subsection{Каталог с нечётким поиском}

\begin{lstlisting}[language=Python, caption={Поиск по каталогу}]
@router.get("/products")
def get_products(q: str = None, category: str = None,
                 min_price: float = None, max_price: float = None,
                 db: Session = Depends(get_db)):
    query = db.query(Product).filter(Product.is_active == True)
    if q:
        for token in q.strip().split():
            query = query.filter(
                Product.name.ilike(f"%{token}%") |
                Product.description.ilike(f"%{token}%")
            )
    if category:
        query = query.filter(Product.category == category)
    if min_price:
        query = query.filter(Product.price >= min_price)
    if max_price:
        query = query.filter(Product.price <= max_price)
    return query.all()
\end{lstlisting}

\subsection{Программа лояльности и заказы}

\begin{lstlisting}[language=Python, caption={Тиры лояльности}]
LOYALTY_TIERS = [
    (10_000_000, 0.50, "Легенда"),
    (5_000_000,  0.25, "Бриллиант"),
    (2_000_000,  0.17, "Изумруд"),
    (1_000_000,  0.13, "Сапфир"),
    (600_000,    0.10, "Платина"),
    (300_000,    0.07, "Золото"),
    (150_000,    0.05, "Серебро"),
    (50_000,     0.03, "Бронза"),
    (0,          0.00, "Новичок"),
]

def get_tier(total_spent: float):
    for threshold, rate, name in LOYALTY_TIERS:
        if total_spent >= threshold:
            return name, rate
    return "Новичок", 0.0
\end{lstlisting}

\subsection{Многоязычный интерфейс (i18n)}

\begin{lstlisting}[language=Java, caption={Конфигурация i18next}]
i18n.use(initReactI18next).init({
  resources: {
    ru: { translation: ru },
    kz: { translation: kz },
    en: { translation: en }
  },
  lng: localStorage.getItem('lang') || 'ru',
  fallbackLng: 'ru',
  interpolation: { escapeValue: false }
});
\end{lstlisting}

\subsection{Защита маршрутов}

\begin{lstlisting}[language=Python, caption={Проверка прав администратора}]
def get_admin_user(
    current_user: User = Depends(get_current_user)
):
    if not current_user.is_admin:
        raise HTTPException(status_code=403,
                            detail="Forbidden")
    return current_user
\end{lstlisting}

\section{Развёртывание}

\textbf{Render (бэкенд):} команда запуска \texttt{uvicorn main:app --host 0.0.0.0 --port \$PORT}. Переменные: \texttt{DATABASE\_URL}, \texttt{SECRET\_KEY}, \texttt{RESEND\_API\_KEY}, \texttt{FRONTEND\_URL}.

\textbf{Vercel (фронтенд):} \texttt{vercel.json} перенаправляет все запросы на \texttt{index.html} для SPA. Переменная \texttt{VITE\_API\_URL} указывает на Render-бэкенд.

% ============================================================
\chapter{ТЕСТИРОВАНИЕ И ОЦЕНКА КАЧЕСТВА}

\section{Методология}

\begin{itemize}
  \item \textbf{Модульное} --- функции хэширования, токенов, алгоритм лояльности
  \item \textbf{Интеграционное} --- API-эндпоинты + база данных
  \item \textbf{Функциональное} --- соответствие всем 24 требованиям
  \item \textbf{UI} --- ручная проверка пользовательских сценариев
\end{itemize}

\section{Тест-кейсы}

\begin{table}[h!]
\centering
\caption{Тест-кейсы аутентификации}
\begin{tabular}{clll}
\toprule
\textbf{№} & \textbf{Тест-кейс} & \textbf{Ожидаемый результат} & \textbf{Статус} \\
\midrule
1 & Регистрация нового пользователя  & 201, письмо отправлено & Passed \\
2 & Регистрация с существующим email & 400                    & Passed \\
3 & Вход с верными данными           & 200, JWT               & Passed \\
4 & Вход с неверным паролем          & 401                    & Passed \\
5 & Сброс пароля                     & 200, письмо отправлено & Passed \\
6 & Сброс --- неизвестный email      & 404                    & Passed \\
\bottomrule
\end{tabular}
\end{table}

\begin{table}[h!]
\centering
\caption{Тест-кейсы лояльности}
\begin{tabular}{clrll}
\toprule
\textbf{№} & \textbf{Сценарий} & \textbf{total\_spent} & \textbf{Уровень} & \textbf{Кешбэк} \\
\midrule
1 & Новый пользователь  &           0 & Новичок   &  0\% \\
2 & После первых покупок &    75\,000 & Бронза    &  3\% \\
3 & Активный покупатель & 350\,000   & Золото    &  7\% \\
4 & VIP-клиент          & 1\,200\,000 & Сапфир   & 13\% \\
5 & Топ-клиент          & 12\,000\,000 & Легенда & 50\% \\
\bottomrule
\end{tabular}
\end{table}

\section{Нагрузочное тестирование}

\begin{table}[h!]
\centering
\caption{Результаты нагрузочного тестирования (50 пользователей)}
\begin{tabular}{lrrr}
\toprule
\textbf{Эндпоинт} & \textbf{Avg RT (мс)} & \textbf{p95 (мс)} & \textbf{RPS} \\
\midrule
GET /api/products    & 142 & 287 & 180 \\
POST /api/auth/token & 234 & 421 & 110 \\
GET /api/orders      & 189 & 334 & 145 \\
\bottomrule
\end{tabular}
\end{table}

Все показатели соответствуют NFR-01 ($\leq$ 500 мс для p95).

\section{Оценка безопасности (OWASP Top 10)}

\begin{table}[h!]
\centering
\caption{Результаты оценки безопасности}
\begin{tabular}{lll}
\toprule
\textbf{Уязвимость} & \textbf{Меры защиты} & \textbf{Статус} \\
\midrule
SQL Injection        & ORM-параметризация (SQLAlchemy) & Защищено \\
XSS                  & React экранирует JSX            & Защищено \\
CSRF                 & SPA + JWT (не cookies)          & Н/А \\
Broken Auth          & bcrypt + JWT exp + email verify & Защищено \\
Sensitive Data       & HTTPS/TLS + хэши паролей        & Защищено \\
Broken Access Control & Role-based guards               & Защищено \\
\bottomrule
\end{tabular}
\end{table}

% ============================================================
\chapter{ЭКОНОМИЧЕСКОЕ ОБОСНОВАНИЕ}

\section{Затраты на разработку}

\begin{table}[h!]
\centering
\caption{Трудозатраты по модулям}
\begin{tabular}{lr}
\toprule
\textbf{Модуль} & \textbf{Трудозатраты (ч)} \\
\midrule
Проектирование архитектуры   & 20 \\
Аутентификация и безопасность & 30 \\
Каталог товаров и поиск      & 25 \\
Корзина и оформление заказов & 35 \\
Программа лояльности         & 20 \\
Многоязычный интерфейс       & 15 \\
Модуль доставки              & 20 \\
Административная панель      & 40 \\
Портал сотрудников           & 20 \\
Тестирование и отладка       & 30 \\
Деплой и конфигурация        & 15 \\
\midrule
\textbf{Итого}               & \textbf{270} \\
\bottomrule
\end{tabular}
\end{table}

При ставке 5\,000 тг/час:
\[\text{Стоимость} = 270 \times 5\,000 = 1\,350\,000 \text{ тг}\]

\section{Операционные расходы}

\begin{table}[h!]
\centering
\caption{Ежемесячные расходы}
\begin{tabular}{lr}
\toprule
\textbf{Статья} & \textbf{Стоимость/мес} \\
\midrule
Vercel (Hobby)                & 0 тг \\
Render Web Service            & $\approx$3\,500 тг \\
Render PostgreSQL             & $\approx$4\,400 тг \\
Resend (до 3\,000 писем/мес)  & 0 тг \\
\midrule
\textbf{Итого}                & $\approx$\textbf{7\,900 тг} \\
\bottomrule
\end{tabular}
\end{table}

\section{Сравнение с альтернативами}

\begin{table}[h!]
\centering
\caption{Сравнение стоимости владения}
\begin{tabular}{lrrl}
\toprule
\textbf{Решение} & \textbf{Внедрение} & \textbf{/мес} & \textbf{Кастомизация} \\
\midrule
TechStore (данная работа) & 1\,350\,000 тг & 7\,900 тг  & Полная \\
Shopify Business          & 0              & 50\,000 тг & Ограниченная \\
Magento Enterprise        & 5\,000\,000 тг & 200\,000 тг & Полная \\
Kaspi магазин             & 0              & 10--20\% продаж & Минимальная \\
\bottomrule
\end{tabular}
\end{table}

При обороте 1\,000\,000 тг/мес окупаемость относительно Kaspi достигается за \textbf{1,5 месяца}.

% ============================================================
\chapter*{ЗАКЛЮЧЕНИЕ}
\addcontentsline{toc}{chapter}{ЗАКЛЮЧЕНИЕ}

В рамках диссертации разработана многофункциональная информационная система электронной торговли «TechStore», адаптированная к условиям казахстанского рынка.

\textbf{Достигнутые результаты:}
\begin{enumerate}
  \item Проведён сравнительный анализ платформ; выявлены ограничения для локального рынка.
  \item Спроектирована реляционная модель на 10 таблицах.
  \item Реализован REST API (60+ эндпоинтов) с JWT и ролевым доступом.
  \item Разработан SPA-фронтенд на React с поддержкой трёх языков.
  \item Реализована девятиуровневая программа лояльности (0\%--50\% кешбэка).
  \item Создан модуль доставки для 55 городов с O(1) поиском.
  \item Развёрнута облачная инфраструктура с CI/CD (Vercel + Render).
  \item Все 24 требования выполнены; p95 времени ответа --- 287 мс.
  \item Операционные расходы --- 7\,900 тг/мес; окупаемость --- 1--2 месяца.
\end{enumerate}

\textbf{Направления развития:} интеграция с Kaspi Pay / Freedom Pay, рекомендательная система, мобильное приложение (React Native), аналитический модуль.

% ============================================================
\begin{thebibliography}{40}
\addcontentsline{toc}{chapter}{СПИСОК ИСПОЛЬЗОВАННЫХ ИСТОЧНИКОВ}

\bibitem{turban2018} Turban E. et al. \textit{Electronic Commerce 2018}. Springer, 2018.
\bibitem{shopify2024} Shopify Inc. \textit{Shopify Documentation}. docs.shopify.com, 2024.
\bibitem{woocommerce2024} WooCommerce. \textit{Developer Documentation}. woocommerce.com, 2024.
\bibitem{adobe2024} Adobe Commerce. \textit{Magento Architecture Guide}. developer.adobe.com, 2024.
\bibitem{kaspi2024} Kaspi.kz. \textit{Seller Information Portal}. kaspi.kz/seller, 2024.
\bibitem{react2024} Facebook Inc. \textit{React Documentation}. react.dev, 2024.
\bibitem{garrett2022} Garrett J. \textit{SPA Performance: React vs MPA}. Web Performance Weekly, 2022.
\bibitem{fastapi2024} Ramírez S. \textit{FastAPI Documentation: Benchmarks}. fastapi.tiangolo.com, 2024.
\bibitem{postgresql2024} PostgreSQL Group. \textit{PostgreSQL 15 Documentation}. postgresql.org, 2024.
\bibitem{rfc7519} Jones M. et al. \textit{RFC 7519: JSON Web Token}. IETF, 2015.
\bibitem{reichheld2000} Reichheld F., Schefter P. \textit{E-Loyalty}. Harvard Business Review, 2000.
\bibitem{stat2023} Агентство РК по статистике. \textit{Языковая политика РК}. stat.gov.kz, 2023.
\bibitem{12factor} Wiggins A. \textit{The Twelve-Factor App}. 12factor.net, 2017.
\bibitem{fielding2000} Fielding R. \textit{Architectural Styles ... Network-based Software}. UC Irvine, 2000.
\bibitem{martin2017} Martin R. \textit{Clean Architecture}. Prentice Hall, 2017.
\bibitem{fowler2002} Fowler M. \textit{Patterns of Enterprise Application Architecture}. Addison-Wesley, 2002.
\bibitem{newman2021} Newman S. \textit{Building Microservices}. O'Reilly, 2021.
\bibitem{devops2016} Kim G. et al. \textit{The DevOps Handbook}. IT Revolution, 2016.
\bibitem{vercel2024} Vercel Inc. \textit{Vercel Documentation}. vercel.com/docs, 2024.
\bibitem{render2024} Render Inc. \textit{Render Documentation}. render.com/docs, 2024.
\bibitem{resend2024} Resend Inc. \textit{Resend API Reference}. resend.com/docs, 2024.
\bibitem{tailwind2024} TailwindCSS. \textit{Documentation}. tailwindcss.com, 2024.
\bibitem{pydantic2024} Pydantic. \textit{Documentation v2}. docs.pydantic.dev, 2024.
\bibitem{sqlalchemy2024} SQLAlchemy. \textit{2.0 Documentation}. docs.sqlalchemy.org, 2024.
\bibitem{i18next2024} i18next. \textit{Documentation}. i18next.com, 2024.
\bibitem{owasp2021} OWASP Foundation. \textit{OWASP Top Ten 2021}. owasp.org, 2021.
\bibitem{finreg2023} АРРФР РК. \textit{Отчёт о рынке электронных платежей}. finreg.kz, 2023.
\bibitem{digital2023} Мин. цифрового развития РК. \textit{Цифровой Казахстан}. digital.gov.kz, 2023.
\bibitem{bcrypt2024} bcrypt. \textit{Python library documentation}. PyPI, 2024.
\bibitem{docker2024} Docker Inc. \textit{Docker Compose Documentation}. docs.docker.com, 2024.
\bibitem{percival2020} Percival H., Gregory B. \textit{Architecture Patterns with Python}. O'Reilly, 2020.
\bibitem{wieruch2023} Wieruch R. \textit{The Road to React}. Self-published, 2023.
\bibitem{grinberg2018} Grinberg M. \textit{Flask Web Development}. O'Reilly, 2018.
\bibitem{chodorow2019} Chodorow K. \textit{MongoDB: The Definitive Guide}. O'Reilly, 2019.
\bibitem{hunt2019} Hunt A., Thomas D. \textit{The Pragmatic Programmer}. Addison-Wesley, 2019.
\bibitem{ecommerce2023} Статистика Казахстана. \textit{E-commerce рынок РК 2023}. stat.gov.kz, 2023.
\bibitem{reactrouter2024} React Router. \textit{v6 Documentation}. reactrouter.com, 2024.
\bibitem{vite2024} Vite. \textit{Documentation}. vitejs.dev, 2024.
\bibitem{python2024} Python SF. \textit{Python 3.11 Documentation}. docs.python.org, 2024.
\bibitem{iso25010} ISO. \textit{ISO/IEC 25010:2011 --- SQuaRE}. ISO, 2011.

\end{thebibliography}

% ============================================================
\appendix

\chapter{Структура проекта}

\begin{lstlisting}[language={}, caption={Файловая структура TechStore}]
techstore/
├── frontend/
│   ├── src/
│   │   ├── components/
│   │   ├── pages/
│   │   │   ├── Home.jsx
│   │   │   ├── Products.jsx
│   │   │   ├── Cart.jsx
│   │   │   ├── Checkout.jsx
│   │   │   ├── Orders.jsx
│   │   │   ├── AdminPanel.jsx
│   │   │   └── EmployeePortal.jsx
│   │   ├── context/
│   │   │   ├── AuthContext.jsx
│   │   │   └── CartContext.jsx
│   │   ├── i18n/locales/
│   │   │   ├── ru.json
│   │   │   ├── kz.json
│   │   │   └── en.json
│   │   └── App.jsx
│   └── vite.config.js
├── backend/
│   ├── main.py
│   ├── models.py
│   ├── schemas.py
│   ├── auth.py
│   ├── database.py
│   └── routers/
│       ├── auth.py
│       ├── products.py
│       ├── orders.py
│       ├── admin.py
│       ├── cashback.py
│       ├── employee.py
│       ├── delivery.py
│       └── password_reset.py
└── render.yaml
\end{lstlisting}

\chapter{API-документация}

\begin{longtable}{llll}
\toprule
\textbf{Метод} & \textbf{Путь} & \textbf{Авторизация} & \textbf{Описание} \\
\midrule
\endhead
POST   & /api/auth/register              & Нет         & Регистрация \\
POST   & /api/auth/token                 & Нет         & Вход \\
GET    & /api/auth/me                    & JWT         & Профиль \\
GET    & /api/products                   & Нет         & Список товаров \\
GET    & /api/products/\{id\}            & Нет         & Детали товара \\
POST   & /api/orders                     & JWT         & Создать заказ \\
GET    & /api/orders                     & JWT         & Мои заказы \\
GET    & /api/cashback/balance           & JWT         & Баланс кешбэка \\
GET    & /api/cashback/tier              & JWT         & Уровень лояльности \\
GET    & /api/delivery/city/\{city\}     & Нет         & Данные доставки \\
GET    & /api/admin/stats                & JWT (admin) & Статистика \\
POST   & /api/admin/products             & JWT (admin) & Создать товар \\
PUT    & /api/admin/products/\{id\}      & JWT (admin) & Обновить товар \\
DELETE & /api/admin/products/\{id\}      & JWT (admin) & Удалить товар \\
GET    & /api/admin/orders               & JWT (admin) & Все заказы \\
GET    & /api/admin/users                & JWT (admin) & Все пользователи \\
GET    & /api/employee/orders            & JWT (emp.)  & Заказы филиала \\
PATCH  & /api/employee/orders/\{id\}/status & JWT (emp.) & Статус заказа \\
GET    & /api/branches                   & Нет         & Список филиалов \\
\bottomrule
\end{longtable}

\vfill
\noindent\rule{\textwidth}{0.4pt}
\textit{Диссертация выполнена самостоятельно. Все используемые источники помечены ссылками.}

\vspace{1cm}
\noindent\textbf{[ФИО магистранта]} \hfill \underline{\hspace{5cm}} \quad «\underline{\quad}» \underline{\hspace{2cm}} 2026~г.

\vspace{0.5cm}
\noindent\textbf{Научный руководитель} \hfill \underline{\hspace{5cm}} \quad «\underline{\quad}» \underline{\hspace{2cm}} 2026~г.

\end{document}
